\documentclass[10pt]{article}
\usepackage[T1]{fontenc}
\usepackage[utf8]{inputenc}
\usepackage[french]{babel}
\usepackage{amsmath,amssymb,mathrsfs,mathpazo}
\usepackage{geometry}
\geometry{a4paper,margin=2cm}
\usepackage{graphicx,array,tabularx,tabularray,colortbl,multirow,multicol,booktabs,enumitem,pifont,ulem,fancybox,diagbox,cancel}
\usepackage{pstricks,pstricks-add,pst-plot,pst-tree,pst-node,pst-text,pst-func,pst-math,pst-bspline,pst-3dplot,pst-eucl,pst-solides3d}
\let\pscalculate\undefined
\usepackage{pst-bezier}
\usepackage{tikz,pgfplots,tkz-tab}
\usetikzlibrary{arrows,arrows.meta,patterns,calc,positioning,trees,decorations.markings,shapes.geometric,shapes.misc,shapes.arrows,intersections,decorations.pathmorphing,decorations.pathreplacing,backgrounds,shadows,babel}
\pgfplotsset{compat=1.18}
\usepgfplotslibrary{fillbetween}
\usepackage[np]{numprint}
\usepackage{listings,makecell,pas-tableur,twoopt}
\let\barre\undefined
\usepackage[pstricks]{tablvar}
\definecolor{darkgreen}{RGB}{0,100,0}
\providecommand{\red}{\color{red}}
\providecommand{\blue}{\color{blue}}
\providecommand{\green}{\color{green}}
\providecommand{\black}{\color{black}}
\providecommand{\white}{\color{white}}
\providecommand{\gray}{\color{gray}}
\providecommand{\cyan}{\color{cyan}}
\providecommand{\magenta}{\color{magenta}}
\providecommand{\yellow}{\color{yellow}}
\providecommand{\violet}{\color{violet}}
\providecommand{\orange}{\color{orange}}
\providecommand{\euro}{\texteuro}
\providecommand{\EUR}{\texteuro}
\providecommand{\ieme}{\textsuperscript{e}}
\providecommand{\degres}{\ensuremath{^\circ}}
\providecommand{\degree}{\ensuremath{^\circ}}
\providecommand{\vv}[1]{\overrightarrow{#1}}
\providecommand{\Coord}[1]{#1}
\providecommand{\decosix}{}
\providecommand{\starredbullet}{\ensuremath{\bullet}}
\pagestyle{empty}
\setlength{\parindent}{0pt}
\usepackage{adjustbox}\DeclareUnicodeCharacter{2236}{:}\DeclareUnicodeCharacter{2190}{\ensuremath{\leftarrow}}
\providecommand{\R}{\mathbb{R}}
\providecommand{\N}{\mathbb{N}}
\providecommand{\Z}{\mathbb{Z}}
\providecommand{\Q}{\mathbb{Q}}
\providecommand{\D}{\mathbb{D}}
\providecommand{\C}{\mathbb{C}}
\providecommand{\e}{{\mathrm{e}}}
\providecommand{\vect}[1]{\overrightarrow{#1}}
\providecommand{\vectt}[1]{\overrightarrow{\mathrm{#1}}}
\providecommand{\vvt}[1]{\overrightarrow{\mathrm{#1}}}
\providecommand{\np}[2][]{{#2~#1}}
\providecommand{\vv}[1]{{\overrightarrow{#1}}}
\renewcommand{\d}{{\mathrm{d}}}
\providecommand{\degres}{{^\circ}}
\providecommand{\degre}{{^\circ}}
\providecommand{\blue}{\color{blue}}
\providecommand{\red}{\color{red}}
\providecommand{\enskip}{\quad}
\providecommand{\llbracket}{\lbrack\!\lbrack}
\providecommand{\rrbracket}{\rbrack\!\rbrack}
\providecommand{\euro}{{\text{€}}}
\providecommand{\no}{{\text{n°}}}
\providecommand{\interval}[2]{{[#1;#2]}}
\providecommand{\intervFO}[2]{{[#1;#2[}}
\providecommand{\dotfill}{\ldots}

\providecommand{\ds}{\displaystyle}
\providecommand{\textding}[1]{#1}
\providecommand{\Oij}{\(\mathrm{O};\vec{\imath},\vec{\jmath}\)}
\providecommand{\Oijk}{\(\mathrm{O};\vec{\imath},\vec{\jmath},\vec{k}\)}
\providecommand{\Ouv}{\(\mathrm{O};\vec{u},\vec{v}\)}
\providecommand{\intervFF}[2]{\ensuremath{[#1;#2]}}
\renewcommand{\euro}{\ifmmode\text{\texteuro}\else\texteuro\fi}
\DeclareUnicodeCharacter{00B0}{\ifmmode^{\circ}\else\textdegree\fi}
\DeclareUnicodeCharacter{20AC}{\ifmmode\text{\texteuro}\else\texteuro\fi}
% Énoncés : Éducation nationale. Transcription et corrigés : APMEP, Denis Vergès et contributeurs. Découpage et classement : CoopMaths, Sébastien Lozano, Éric Elter et contributeurs. Adaptation de la présentation : Maths Explorer.
% CC BY-SA : https://coopmaths.fr/www/about/
% https://coopmaths.fr/alea/static/eam/2026/tex/eam_technologique_2026_06_polynesie_3_cor.tex
\begin{document}
\noindent\textit{APMEP / CoopMaths. CC BY-SA. Présentation adaptée par Maths Explorer.}\par\medskip
\begin{enumerate}
	\item On développe l'expression factorisée proposée :

	$(0,5x + 1,5)(-x+5) = 0,5x \times (-x) + 0,5x \times 5 + 1,5 \times (-x) + 1,5 \times 5 $

	On réduit les produits :

	$(0,5x + 1,5)(-x+5) = -0,5x^2 + 2,5x - 1,5x + 7,5 $

	On réduit les sommes :

	$(0,5x + 1,5)(-x+5) = -0,5x^2 + x + 7,5 $

	On obtient bien l'expression de $f(x)$ donnée dans le sujet.

	Pour tout $x$ appartenant à l'intervalle $[-4 ; 6]$,
	\[
	f(x)=(0,5x + 1,5)(-x+5) = -0,5x^2 +  x + 7,5.
	\]

	\item Le nombre dérivé en 3 est la pente de la tangente à la courbe au point d'abscisse 3.

	Par lecture graphique, \quad $f'(3) = \dfrac{4 - 6}{4 - 3} = \dfrac{-2}{1} = -2$.

	En donnant l'expression de la fonction dérivée : \quad $f'(x) = -0,5 \times (2x) + 1 = -x+1$.

	On détermine \quad $f'(3) = -3  + 1 = -2$.


	\item \begin{enumerate}
		\item Graphiquement, les points de la courbe d'ordonnée inférieure ou égal à 0 sont les points d'abscisse appartenant à $\left[-4 ; -3\right]\cup\left[5 ; 6\right]$.

		Algébriquement, on utilise la forme factorisée pour déterminer le signe.

		$\begin{aligned}[t]
			0,5x + 1,5 = 0
			&\iff 0,5x = -1,5 \\
			&\iff x = -3
		\end{aligned}$
		\quad et \quad
		$\begin{aligned}[t]
			0,5x + 1,5 > 0
			&\iff 0,5x > -1,5 \\
			&\iff x > -3
		\end{aligned}$

		$-x+5 = 0 \iff x = 5$\quad  et  \quad $\begin{aligned}[t]
			-x+5 > 0 &\iff 5 > x \\&\iff x < 5.
		\end{aligned}$

	Sur l'intervalle $[-3 ; 5]$,\quad $0,5x + 1,5 \geqslant 0$ \quad  et \quad $-x+5 \geqslant 0$

	donc le produit \quad $(0,5x + 1,5)(-x+5)$ \quad est positif.

	\textbf{L'affirmation~a est fausse.}

	\item Graphiquement, on voit que la fonction est décroissante sur l'intervalle $[1 ; 4]$, donc le nombre dérivé $f'(x)$ est négatif sur l'intervalle $[1 ; 4]$.

	On a vu que \quad $f'(3) = -2$ \quad donc $f'(x)$ n'est pas positif sur l'intervalle $[1 ; 4]$.

	\textbf{L'affirmation~b est fausse.}
	\end{enumerate}
\end{enumerate}

\end{document}
